\chapter{Introduction}

Vision Transformers (ViTs) have become the preferred architecture in computer vision tasks, due to their ability to capture Long-Range dependencies.[2] While ViTs achieve state-of-the-art accuracy, their computation and memory requirements remain significantly higher than traditional \textbf{Convolutional Neural Networks[CNNs].} The primary reason is the self-attention mechanism, whose complexity scales quadratically ($O(n^2)$) with the number of image patches (tokens). As Image resolution increases, the number of patches also increases proportionally, resulting in \textbf{high Latency and Resource usage}. This makes real time deployment on resource constrained environments like FPGAs a big challenge.

\section{Background}

The transition from Convolutional Neural Networks (CNNs) to Vision Transformers (ViTs) has marked a new era in computer vision [4] enabling superior representation learning and improved accuracy across classification, object detection, and segmentation tasks. However, this advancement comes with substantial computational overhead. For vision applications targeting embedded or hardware-accelerated platforms, such as Field-Programmable Gate Arrays (FPGAs), the unoptimized ViT architecture poses a major bottleneck.


\section{Motivation}

Although Vision Transformers offer substantial advantages over CNNs, their computational inefficiency limits their applicability in edge computing and FPGA-based accelerators. The core motivations of this work are summarized below:

\begin{enumerate}

    \item \textbf{Addressing ViT’s Computational Bottleneck:}
    The self-attention mechanism scales quadratically with the number of tokens. As image resolutions increase, this computation becomes prohibitively expensive for real-time or embedded applications without architectural optimization.

    \item \textbf{Need for Efficient Token Reduction:}
    Many tokens—especially those belonging to low-information background regions—contribute minimally to the final prediction. Removing such tokens can significantly reduce compute and memory requirements without degrading accuracy.


    \item \textbf{Direct C-to-RTL Conversion Using HLS:}
    High-Level Synthesis (HLS) enables direct translation of C/C++ code into synthesizable RTL, eliminating the need for manually writing complex Verilog/VHDL descriptions. This drastically reduces design time and allows rapid exploration of architectural optimizations such as pipelining, array partitioning, and loop unrolling.

    \item \textbf{Bridging the Algorithm–Hardware Gap:}
    Existing token pruning techniques are largely optimized for GPUs. There is limited work on adapting these methods for hardware-aware design flows. Integrating pruning directly into the HLS-based ViT pipeline enables efficient deployment on resource-constrained accelerators.

\end{enumerate}

These motivations drive an \textbf{algorithm–hardware co-design} approach where token pruning and hardware optimization work together to enable efficient, real-time ViT inference on FPGA platforms.



\section{Project Objectives}

The primary objective of this project is to design, implement, and validate a hardware accelerator for the Vision Transformer (ViT) model using C based High-Level Synthesis (HLS).

The specific objectives are as follows:
\begin{enumerate}
    \item \textbf{Develop a Parameterized ViT Core in HLS:}  Develop modular and reusable HLS implementations of the fundamental ViT components—including patch embedding, Layer Normalization, Multi-Head Self-Attention (MHSA), and the MLP block.
    \item \textbf{Integrate Dynamic Token Pruning:} Implement an Attention-based Token Selection (ATS) to prune less important tokens/patches dynamically based on their attention scores.
    \item \textbf{Validate Accelerator Performance and Accuracy:} Evaluate the HLS based token pruned ViT accelerator against base line ViT-Tiny, ViT-Small, and ViT-Base models  to compare their accuracies, througput and resource usage .
    \item \textbf{Analyze HLS Optimization Strategies:} To apply and analyze the key HLS pragmas (\texttt{PIPELINE}, \texttt{ARRAY\_PARTITION}, \texttt{UNROLL}) required to transform the algorithmic C or C++ code into a parallel, high-throughput architecture.
\end{enumerate}

\section{Contributions}

This thesis makes the following key contributions:

\begin{enumerate}
    \item \textbf{HLS-Based Vision Transformer Architecture:}
    A complete and parameterized ViT-Tiny architecture is implemented in C/C++ using High-Level Synthesis. The design includes patch embedding, LayerNorm, Multi-Head Self-Attention (MHSA), MLP layers, and residual connections.

    \item \textbf{Hardware-Friendly Token Pruning Module:}
    An Attention-based Token Selection (ATS) mechanism is integrated within the ViT pipeline, enabling dynamic reduction of tokens across layers while preserving classification accuracy.

    \item \textbf{Optimized Self-Attention Engine for FPGA:}
    A fully pipelined and parallelized MHSA engine is developed using HLS pragmas such as \texttt{PIPELINE}, \texttt{UNROLL}, and \texttt{ARRAY\_PARTITION}, resulting in significant reductions in latency.

    \item \textbf{Comprehensive Experimental Evaluation:}
    The proposed accelerator is validated against baseline ViT-Tiny, ViT-Small, and ViT-Base models. Comparisons include accuracy, throughput, and FPGA resource utilization, demonstrating the benefits of token pruning for ViT acceleration.
\end{enumerate}

\section{Organization of The Report}

This chapter provided a background for the topics covered in this report, describing the computational challenges of Vision Transformers and the motivation for FPGA acceleration. The rest of the chapters are organized as follows:

\textbf{Chapter 2} provides a fundamental introduction to Vision Transformers, and how they differ from CNNs.
\textbf{Chapter 3} provides a review of prior works, specifically focusing on existing token pruning strategies and hardware accelerators.
\textbf{Chapter 4} explains the Attention mechanism, along with the mathematics behind the $O(n^2)$ bottleneck and how attention scores can be used for pruning thus reducing number of tokens layer by layer.
\textbf{Chapter 5} discusses the HLS implementation, explaining pragmas and mathematical approximations for creating hardware circuit.
\textbf{Chapter 6} presents the Experimental Results, including Accuracy, latency and resource utilization.
Finally, \textbf{Chapter 7} concludes the report and discusses future scope of this project.